\documentclass[12pt]{article}
\usepackage[utf8]{inputenc}
\usepackage[T1]{fontenc}
\usepackage[margin=1in]{geometry}
\usepackage{setspace}
\usepackage{xcolor}
\usepackage[colorlinks=true, urlcolor=blue!60!black, linkcolor=blue!60!black]{hyperref}
\setlength{\parindent}{0em}
\setlength{\parskip}{0.6em}
\renewcommand{\baselinestretch}{1.0}
\begin{document}
\begin{center}
{\LARGE \textit{Grutter v. Bollinger}, 539 U.S. 306 (2003)\par}
\vspace{0.5em}
{\large Case Brief --- Two-Page Edition\par}
\end{center}

\section{Memory Jogger}
Student body diversity is a compelling interest, so Michigan Law School's holistic, race-conscious admissions survived strict scrutiny 5--4---until \textit{SFFA v. Harvard} (2023) (\href{https://www.courtlistener.com/opinion/130155/grutter-v-bollinger/}{CourtListener}).

\section{Detailed Case Facts}
Michigan Law School's 1992 policy, modeled on Justice Powell's \textit{Bakke} opinion, reviewed each file holistically and sought a ``critical mass'' of underrepresented minority students. Barbara Grutter, a white applicant with a 3.8 GPA and a 161 LSAT, was rejected and sued under the Equal Protection Clause, Title VI, and \S~1981.

\section{Procedural History}
The district court held the policy unconstitutional and enjoined it; the en banc Sixth Circuit reversed 5--4. The Supreme Court affirmed, the same day it struck down Michigan's undergraduate point system in \textit{Gratz}.

\section{Judicial Votes}
Affirmed 5--4. O'Connor wrote for Stevens, Souter, Ginsburg, and Breyer; Ginsburg concurred; Scalia and Thomas concurred in part and dissented in part; Rehnquist (with Scalia, Kennedy, and Thomas) and Kennedy dissented.

\section{Holding}
Diversity is a compelling interest in higher education. A program is narrowly tailored if it uses race as a flexible ``plus'' in individualized review, seriously considers race-neutral alternatives, does not unduly burden nonminority applicants, and is limited in duration.

\section{Analysis of Opinions}
\textbf{Majority:} The Court deferred to the school's educational judgment; diversity's benefits were ``not theoretical but real,'' and critical mass was not a quota. \textbf{Ginsburg:} ``one may hope, but not firmly forecast'' an end in 25 years. \textbf{Thomas:} only ``pressing public necessity'' justifies racial classifications. \textbf{Rehnquist:} admissions tracked applicant shares, proving racial balancing. \textbf{Kennedy:} ``Deference is antithetical to strict scrutiny, not consistent with it.''

\section{Examples: Future Applications}
\textbf{Same-Side:} holistic medical school review with race as one plus factor (now invalid after \textit{SFFA}). \textbf{Opposite-Side:} a 15-point racial bonus, like \textit{Gratz}. \textbf{Fence-Sitter:} a ZIP-code preference adopted to preserve racial diversity.

\section{Critique}
Conservatives said deference diluted strict scrutiny, and Richard Sander's contested ``mismatch'' thesis claimed preferences hurt beneficiaries. From the left, Derrick Bell argued that diversity was a ``distraction'' from remedying racial subordination.

\section{Key Quotations}
\begin{quote}
``We expect that 25 years from now, the use of racial preferences will no longer be necessary to further the interest approved today.''
\end{quote}
\noindent --- Justice O'Connor, Opinion of the Court, 539 U.S. at 343.

\end{document}
